\documentclass[12pt]{article}
\usepackage[utf8]{inputenc}
\usepackage[T1]{fontenc}
\usepackage{hyperref}
\title{Berawan, East Language Resources for K-12 Curriculum Development}
\author{Compiled by Claude Sonnet 5.5 (claude-sonnet-5-5) for Andr\'as Kornai}
\date{}
\begin{document}
\maketitle

\section{Introduction}
East Berawan (ISO 639-3 zbe; Glottocode east2486) is an endangered Austronesian language of Sarawak, Malaysia. In Burkhardt's reconstruction of Proto-Berawan, East Berawan corresponds to the Long Terawan lect, while Batu Belah, Long Teru and Long Jegan form a separate western group. This is how the sources describe the classification; the exact mapping of lects to ISO codes was not confirmed in the searches. Little material specific to East Berawan was found, and most resources concern Berawan as a whole.

\section{Primary Educational Resources}
No textbooks or primers were found. A unified orthography for the four Berawan lects has been published and could serve as a basis.
\begin{itemize}
\item One Orthography, Four Lects: The Unified Berawan Orthography. \url{https://publisher.unimas.my/ojs/index.php/BJK/article/view/2411}
\item Full text of the same article. \url{https://publisher.unimas.my/ojs/index.php/BJK/article/view/2411/1012}
\end{itemize}

\section{Linguistic and Reference Materials}
\begin{itemize}
\item Wikipedia article on the Berawan language. \url{https://en.wikipedia.org/wiki/Berawan_language}
\item Burkhardt, paper on Terawan Berawan diphthongs. \url{https://sil-philippines-languages.org/ical/papers/burkhardt-terawan_berawan_dipthong.pdf}
\item Robert Blust, Berawan recordings including Long Terawan (1971). \url{https://researchdata.edu.au/rb1-003-robert-blust-recording/1577760}
\item PARADISEC catalogue item RB1/003. \url{https://catalog.paradisec.org.au/repository/RB1/003}
\item Scholarspace (University of Hawai\textquoteleft i) record that appeared in the search for Berawan; its content was not examined. \url{https://scholarspace.manoa.hawaii.edu/items/d68956e1-55ae-4d57-94a6-95ffb833d0cc/full}
\item Sarawak Museum web page on Berawan. \url{https://museum.sarawak.gov.my/web/subpage/webpage_view/847}
\end{itemize}

\section{Religious and Community Resources}
No East Berawan specific religious resources were found in the searches. Global Recordings and Joshua Project entries found were for Central Berawan only.

\section{Digital and Audio Resources}
The Blust recordings above are the main audio found.

\section{Curriculum Development Context}
East Berawan is still spoken by all adults but not all young people, and is not known to be taught in schools. A curriculum would combine revitalization with documentation, starting from the unified orthography and with community involvement. Existing school materials were not found.
\end{document}